% R4 — in-distribution positive control on monocytes (draft 2026-09-07)
\subsection{The in-distribution positive control returns no separation either}
\label{ssec:r4_monocyte}

A reviewer could dismiss the platelet results as an artifact of evaluating too
far outside the models' intended use. We therefore reran the identical pipeline
on a cell type the models were trained on: monocytes from the same source
cohorts (Ren and Stephenson COVID-19; five Qiu sepsis patients; 148
patients\footnote{One Ren survivor patient (S-S072) is not represented in
the monocyte cell set, so the monocyte cohort contains 148 rather than 149
patients and the disease-matched COVID-only subset 143 rather than 144;
fatal counts are unaffected.}, 19 fatal; Section~\ref{ssec:methods_baselines}). Monocytes are among the most
abundant immune cells in the atlases these models were pretrained on, so this
arm isolates cell-type familiarity as the single changed variable.

The positive control returned no separation at all
(Table~\ref{tab:monocyte} and Figure~\ref{fig:indist_vs_ood}). Under
patient-level paired bootstrap, no foundation model significantly outperformed
the PCA-200 baseline on monocytes (\model{STATE} $+0.043$, 95\% CI [$-0.026$,
0.120], $p = 0.23$; \model{Geneformer} $+0.022$, $p = 0.57$; \model{scPRINT-2}
$+0.018$, $p = 0.68$; \model{UCE} $+0.012$, $p = 0.80$; \model{scGPT} $+0.005$,
$p = 0.94$). Even on a cell type these models were trained on, the from-scratch
baseline is statistically sufficient. One directional pattern is worth
noting: \model{UCE}'s monocyte-versus-platelet difference is the closest to
zero of any method in both cohort pools ($-0.005$ COVID-only, $-0.014$
all-cohort, versus drops of $0.043$--$0.119$ for the other five methods) ---
the least affected by cell type, though it beats the baseline on neither.

One inversion requires careful reading. Every method, including the baseline,
scored lower on monocytes than on platelets (baseline $-0.072$; confirmed on
the disease-matched COVID-only subset, 143 patients, 17 fatal). Because the
from-scratch baseline drops by as much as the foundation models, the inversion
is a property of the task and cohort, not of pretraining: patient-level
mortality from monocytes is simply a harder measurement in these cohorts. That
platelet transcriptomes carry more mortality signal than monocyte
transcriptomes in COVID-19 is consistent with the mechanistic literature
placing platelet-driven immunothrombosis proximal to fatal disease
\citep{manne2020platelet,barrett2021platelets}. The positive control therefore narrows the out-of-distribution
objection in an unexpected direction. Even on a cell type the models were
trained on, the from-scratch baseline is statistically sufficient, and with
the within-COVID primary contrast also returning null under the
winner-corrected test (Section~\ref{ssec:r1_within_covid}), the pipeline
resolved no pretraining advantage over the baseline in any cell of the
design. The platelet results can therefore not be attributed to evaluating
too far outside the models' intended use: the same comparison is null
in-distribution. Equally, the control cannot certify that this pipeline
resolves foundation-model advantages anywhere, and we report it with that
qualification (Section~\ref{ssec:limitations}).
